% ===========================================================================
% exp06_hdbscan.tex - complete code listings for Experiment 6
% Source: notebooks/02_density_based_learning.ipynb (executed)
% ===========================================================================

\begin{lstlisting}[style=mlcode,caption={Core numerical and plotting libraries}]
# ---- Core numerical and plotting libraries ----
import numpy as np
import pandas as pd
import matplotlib.pyplot as plt
import seaborn as sns

# ---- Density-based clustering algorithms ----
from sklearn.cluster import DBSCAN, HDBSCAN
# ---- Nearest-neighbor search (used to pick epsilon) ----
from sklearn.neighbors import NearestNeighbors
# ---- Evaluation metrics ----
from sklearn.metrics import silhouette_score

# Aesthetics setup (consistent look across all notebooks)
plt.style.use('seaborn-v0_8-whitegrid' if 'seaborn-v0_8-whitegrid' in plt.style.available else 'default')
plt.rcParams['font.sans-serif'] = 'DejaVu Sans'
plt.rcParams['figure.figsize'] = (11, 6)
plt.rcParams['figure.dpi'] = 110

np.random.seed(42)
print("Density-based libraries loaded successfully!")
\end{lstlisting}

\begin{lstlisting}[style=mlcode,caption={Load the real USGS earthquake catalogue (global, 2023)}]
# ---- Load the real USGS earthquake catalogue (global, 2023) ----
df = pd.read_csv('../data/earthquakes.csv')
print(f"Full catalogue (magnitude >= 4.0): {len(df):,} events")
print(df.head(3).to_string(index=False))

# Keep well-recorded moderate+ events: magnitude >= 4.5
quakes = df[df['mag'] >= 4.5].reset_index(drop=True)
print(f"\nAnalysed subset (magnitude >= 4.5): {len(quakes):,} events")

# ---- Convert epicenters to radians [latitude, longitude] ----
# scikit-learn's 'haversine' metric expects radian coordinates and computes
# great-circle distance on the sphere (correct for global data).
X_rad = np.radians(quakes[['latitude', 'longitude']].values)
print("Coordinate range (deg): lat", quakes.latitude.min(), "to", quakes.latitude.max(),
      "| lon", quakes.longitude.min(), "to", quakes.longitude.max())
\end{lstlisting}

\begin{lstlisting}[style=mlcode,caption={Fit HDBSCAN (no global epsilon required)}]
# ---- Fit HDBSCAN (no global epsilon required) ----
# min_cluster_size = smallest group considered a cluster; min_samples controls conservativeness.
hdb = HDBSCAN(min_cluster_size=50, min_samples=10, metric='haversine', copy=False)
hdb_labels = hdb.fit_predict(X_rad)

n_clusters_hdb = len(set(hdb_labels)) - (1 if -1 in hdb_labels else 0)
n_noise_hdb = np.sum(hdb_labels == -1)
print(f"HDBSCAN: {n_clusters_hdb} clusters, {n_noise_hdb:,} noise events ({n_noise_hdb / len(X_rad) * 100:.1f}%)")

# ---- Visualize the HDBSCAN partition ----
plt.figure(figsize=(12, 5.5))
palette = sns.color_palette("tab20", max(n_clusters_hdb, 1))
for k in range(n_clusters_hdb):
    mask = hdb_labels == k
    plt.scatter(quakes.longitude[mask], quakes.latitude[mask], s=6, color=palette[k % 20], alpha=0.75)

noise_mask = hdb_labels == -1
plt.scatter(quakes.longitude[noise_mask], quakes.latitude[noise_mask], s=6, color='black', marker='x', label='Noise (-1)')
plt.title(f"HDBSCAN on Global Earthquakes 2023 ({n_clusters_hdb} clusters, no ε)", fontsize=13, fontweight='bold')
plt.xlabel("Longitude", fontsize=12)
plt.ylabel("Latitude", fontsize=12)
plt.legend(loc='lower left')
plt.tight_layout()
plt.show()
\end{lstlisting}
